\documentclass[12pt,a4paper]{article}
\usepackage[utf8]{utf8}
\usepackage[spanish,es-tabla]{babel}
\usepackage{geometry}
\geometry{top=2.5cm,bottom=2.5cm,left=3cm,right=2.5cm}
\usepackage{graphicx}
\usepackage{booktabs}
\usepackage{tabularx}
\usepackage{hyperref}
\usepackage{listings}
\usepackage{xcolor}

\definecolor{primarycolor}{RGB}{31, 73, 125}
\definecolor{secondarycolor}{RGB}{89, 89, 89}
\definecolor{codebg}{RGB}{245, 245, 245}

\hypersetup{
    colorlinks=true,
    linkcolor=primarycolor,
    citecolor=primarycolor,
    urlcolor=primarycolor
}

\begin{document}

% -----------------------------------------------------------------
% PORTADA DEL MANUAL DE USUARIO
% -----------------------------------------------------------------
\begin{titlepage}
    \centering
    \vspace*{1.5cm}

    {\Huge \bfseries \color{primarycolor} ASISTDOC} \\[0.5cm]
    {\Large \bfseries Sistema Local para la Gestión de Asistencia y Calificaciones Académicas} \\[1.5cm]

    \rule{\linewidth}{0.5mm} \\[1cm]

    {\huge \bfseries MANUAL DE USUARIO PARA EL DOCENTE} \\[2cm]

    \begin{minipage}{0.48\textwidth}
        \begin{flushleft} \large
            \textbf{Manual Operativo} \\[0.2cm]
            Versión 1.0 (Local / Offline)
        \end{flushleft}
    \end{minipage}
    \begin{minipage}{0.48\textwidth}
        \begin{flushright} \large
            \textbf{Público Objetivo:} \\[0.2cm]
            Docentes Universitarios e Instructores
        \end{flushright}
    \end{minipage}

    \vfill

    {\large 2026}
\end{titlepage}

\tableofcontents
\newpage

% -----------------------------------------------------------------
% 1. INTRODUCCIÓN Y BIENVENIDA
% -----------------------------------------------------------------
\section{Introducción}
Bienvenido a \textbf{AsistDoc}, el sistema de escritorio de uso local diseñado especialmente para facilitarle la gestión integral de sus asignaturas, el registro rápido de asistencia, la consolidación de calificaciones y la emisión de reportes académicos oficiales.

AsistDoc funciona de manera 100\% local y privada en su computadora personal. No requiere conexión a internet ni suscripciones a servicios en la nube, garantizando que la información de sus estudiantes permanezca siempre segura y bajo su total control.

% -----------------------------------------------------------------
% 2. REQUISITOS DEL SISTEMA
% -----------------------------------------------------------------
\section{Requisitos del Sistema}
Para utilizar AsistDoc en su equipo únicamente necesita:
\begin{itemize}
    \item \textbf{Sistema Operativo:} Windows 10/11, macOS (10.15 o superior) o Linux.
    \item \textbf{Interpretador de Python:} Python versión 3.10 o superior.
    \item \textbf{Navegador Web:} Google Chrome, Mozilla Firefox, Microsoft Edge o Safari (para visualizar la interfaz de la aplicación).
\end{itemize}

% -----------------------------------------------------------------
% 3. GUÍA DE INSTALACIÓN PASO A PASO
% -----------------------------------------------------------------
\section{Instalación del Sistema}

\subsection{En Sistemas Linux / macOS}
Abra su terminal de comandos e ingrese la siguiente secuencia:
\begin{verbatim}
cd AsistDoc
python3 -m venv .venv
source .venv/bin/activate
pip install -r requirements.txt
\end{verbatim}

\subsection{En Sistemas Windows}
Abra la Consola de Comandos (CMD) o PowerShell e ingrese:
\begin{verbatim}
cd AsistDoc
python -m venv .venv
.venv\Scripts\activate
pip install -r requirements.txt
\end{verbatim}

% -----------------------------------------------------------------
% 4. INICIO Y NAVEGACIÓN PRINCIPAL
% -----------------------------------------------------------------
\section{Inicio y Navegación Principal}

\subsection{Iniciar la Aplicación}
Para iniciar AsistDoc, ejecute en su terminal el comando:
\begin{verbatim}
streamlit run app.py
\end{verbatim}
Inmediatamente se abrirá una ventana en su navegador web predeterminado mostrando el panel principal.

\subsection{Barra Lateral de Navegación}
En el margen izquierdo encontrará la barra lateral con los 6 módulos principales:
\begin{enumerate}
    \item \textbf{Hoy:} Visión general de las clases programadas para el día actual.
    \item \textbf{Cursos y alumnos:} Gestión de asignaturas, horarios e importación de listas.
    \item \textbf{Asistencia:} Toma rápida de asistencia y cierre de sesión.
    \item \textbf{Notas:} Matriz de calificaciones, evaluaciones e Ingreso Rápido.
    \item \textbf{Reportes:} Motor de alertas de riesgo, gráficos y exportaciones Excel/PDF.
    \item \textbf{Configuración:} Copias de seguridad y auditoría de cambios.
\end{enumerate}

% -----------------------------------------------------------------
% 5. GESTIÓN DE CURSOS Y ALUMNOS
% -----------------------------------------------------------------
\section{Gestión de Cursos y Alumnos}

\subsection{Creación de un Curso}
1. Ingrese a \textbf{Cursos y alumnos} $\rightarrow$ Pestaña \textbf{Cursos y Horarios}.\\
2. Complete el formulario con el Nombre del curso (ej. \textit{Redes II}), Sección (ej. \textit{A}), Nota mínima aprobatoria (ej. \textit{11.0}) y Límite de Faltas (ej. \textit{30\%}).\\
3. Presione el botón \textbf{Guardar Curso}.

\subsection{Importación Masiva de Alumnos desde Excel o CSV}
Para matricular a todos los alumnos de un grupo sin escribirlos uno a uno:
1. Diríjase a la pestaña \textbf{Importación Masiva}.\\
2. Prepare un archivo Excel (.xlsx) o CSV con al menos dos columnas obligatorias: \texttt{codigo} y \texttt{nombre}.\\
3. Suba el archivo en el recuadro correspondiente.\\
4. Verifique la vista previa en pantalla y presione \textbf{Confirmar Importación Masiva}.

% -----------------------------------------------------------------
% 6. TOMA DE ASISTENCIA
% -----------------------------------------------------------------
\section{Toma y Control de Asistencia}
1. Seleccione el curso y la fecha correspondiente.\\
2. Si todos los alumnos asistieron, presione el botón \textbf{"Todos presentes"} para marcarlos de un solo clic.\\
3. Si algún alumno llegó tarde o faltó, modifique únicamente su casilla individual. El sistema calculará automáticamente la \textbf{tardanza} si la hora de registro supera la tolerancia configurada.\\
4. Al finalizar la clase, presione \textbf{"Cerrar Asistencia"}. Los alumnos no registrados se marcarán automáticamente como \textit{falta}.

% -----------------------------------------------------------------
% 7. MÓDULO DE NOTAS E INGRESO RÁPIDO
% -----------------------------------------------------------------
\section{Módulo de Calificaciones e Ingreso Rápido}

\subsection{Modo "Ingreso Rápido" (Función Estrella)}
Este modo está diseñado para registrar notas de forma ágil sin usar el ratón:
1. Ingrese a \textbf{Notas} $\rightarrow$ Pestaña \textbf{INGRESO RÁPIDO}.\\
2. Seleccione la evaluación activa (ej. \textit{Práctica 2}).\\
3. En el buscador escriba el apellido o código del alumno (ej. \textit{mamani}). El sistema encuentra coincidencias ignorando mayúsculas y tildes.\\
4. Introduzca la nota (de 0 a 20) o escriba \textbf{NP} si el alumno no se presentó, y presione \textbf{Enter}.\\
5. La nota quedará registrada y el formulario se limpiará listo para el siguiente alumno.

\subsection{Deshacer Último Ingreso}
Si cometió un error de digitación, presione el botón \textbf{"Deshacer Último Ingreso"} para revertir inmediatamente la última calificación guardada.

\subsection{Calculadora "¿Cuánto Necesita?"}
En la pestaña \textbf{"¿Cuánto Necesita?"}, seleccione un estudiante y una evaluación pendiente. El sistema le indicará exactamente qué calificación sobre 20 requiere para alcanzar la nota mínima aprobatoria considerando los pesos del curso.

% -----------------------------------------------------------------
% 8. REPORTES Y EXPORTACIÓN
% -----------------------------------------------------------------
\section{Reportes, Alumnos en Riesgo y Exportación}
1. Diríjase al módulo \textbf{Reportes}.\\
2. En la pestaña \textbf{Alumnos en Riesgo}, consulte la lista de estudiantes clasificados como \textbf{ALTO RIESGO} o \textbf{ATENCIÓN} debido a bajo rendimiento o exceso de inasistencias.\\
3. En la pestaña \textbf{Descargar Informes}, presione \textbf{"Generar Acta Excel"} para obtener la matriz completa con formato profesional de colores (\textit{.xlsx}) o \textbf{"Generar Acta PDF"} para la versión oficial imprimible.

% -----------------------------------------------------------------
% 9. COPIAS DE SEGURIDAD
% -----------------------------------------------------------------
\section{Copias de Seguridad}
Para proteger su información académica:
1. Vaya a \textbf{Configuración} $\rightarrow$ Pestaña \textbf{Copias de Seguridad}.\\
2. Presione \textbf{"Crear Copia de Seguridad Ahora"}. El sistema generará un archivo con fecha y hora en la carpeta \texttt{backups/}.\\
3. Si necesita restaurar una copia previa, selecciónela de la lista y confirme la restauración.

% -----------------------------------------------------------------
% 10. SOLUCIÓN DE PROBLEMAS FRECUENTES
% -----------------------------------------------------------------
\section{Solución de Problemas Frecuentes}

\begin{table}[h!]
\centering
\caption{Guía de Solución de Inconvenientes Frecuentes}
\begin{tabularx}{\textwidth}{X X}
\toprule
\textbf{Problema / Mensaje} & \textbf{Solución Sugerida} \\
\midrule
\textit{"Puerto 8501 ocupado"} & Inicie Streamlit en un puerto alternativo: \texttt{streamlit run app.py --server.port 8502}. \\
\textit{"El archivo debe contener las columnas 'codigo' y 'nombre'"} & Abra su archivo Excel de importación y asegúrese de que la primera fila tenga exactamente las cabeceras \texttt{codigo} y \texttt{nombre}. \\
\textit{"La nota debe estar entre 0.0 y 20.0"} & Verifique el valor ingresado. Si el estudiante no rindió la prueba, escriba la palabra \textbf{NP}. \\
\textit{"Suma de pesos no es 100\%"} & Vaya a Configurar Evaluaciones y ajuste los pesos porcentuales para que su suma alcance exactamente el 100\%. \\
\bottomrule
\end{tabularx}
\end{table}

\end{document}
